\clearpage

\section {Introduction }
\label{sec:intro}

In this technical note, 
the results of a joint 3-flavour oscillation analysis
performed by the {\em VALOR} group
on the combined Run 1+2+3+4 (6.5700$\times$10$^{20}$ POT) dataset are reported.
This dataset was collected in the following run periods:
in January - June 2010 (Run1, 0.323$\times$10$^{20}$ POT),
in November 2010 - March 2011 (Run2, 1.108$\times$10$^{20}$ POT) and,
following the recovery work from the damage of the 2011 Tohoku earthquake \cite{earthquake}, 
in March - June 2012 (Run3, 1.579$\times$10$^{20}$ POT) 
and in October 2012 - 8th May 2013 (Run 4, 3.560$\times$10$^{20}$ POT).

This analysis is an update of the work presented in Ref. \cite{OurJointAnalysis}, based on the 
{\em VALOR}  $\nu_{\mu}$-disappearance analysis presented in Refs. \cite{OurRun1DisappAnalysis}, 
\cite{OurRun12DisappAnalysis}, \cite{OurRun12DisappAnalysis3Flavour} and \cite{OurRun123DisappAnalysis3Flavour},
with improved methods and an enlarged dataset. The baseline of this analysis has been explained in detail in the Run 1+2+3 technical note (Ref. \cite{OurJointAnalysis}); a summary of the updates in the analysis inputs and 
structural changes concerning the fitting method, such as adding together SK efficiencies and FSI+SI systematics,
are summarized in Sec.\ref{sec:changes}. Among these changes, the most remarkable one is that this updated joint analysis
contains the correlations between the SK efficiencies for the single $\mu$-like ring and the single e-like ring event samples that were not available in the previous analysis for Run 1+2+3 dataset.

In the joint oscillation analysis described in this technical note,
the atmospheric squared-mass splitting $|\Delta m^{2}|$,
the mixing angles $\sin^{2}\theta_{23}$, $\sin^{2}\theta_{13}$
and the CP-violating phase $\delta_{CP}$,
are determined together through a simultaneous fit of the reconstructed
energy spectra of both the SuperK single $\mu$-like ring and single e-like ring event samples.
In contrast to the stand-alone $\nu_{e}$-appearance and  $\nu_{\mu}$-disappearance fits,
the joint 3-flavour oscillation fit presented in this work does not have to assume prior values for any of these 
four oscillation parameters and therefore all correlations between them are taken into account.
The 12-sector (solar) oscillation parameters, $\sin^{2}\theta_{12}$ and $|\delta m^{2}|$ are kept fixed 
since their effect was demonstrated to be negligible in \ref{app:solar}.

This analysis is a continuation of the work of the {\em VALOR} $\nu_{\mu}$-disappearance analysis presented in
Refs.\ \cite{OurRun1DisappAnalysis}, \cite{OurRun12DisappAnalysis}, \cite{OurRun12DisappAnalysis3Flavour}
and \cite{OurRun123DisappAnalysis3Flavour}. 
The {\em VALOR} 3-flavour stand-alone $\nu_{\mu}$-disappearance analysis framework was extended 
to include a new 3-flavour $\nu_{e}$-appearance analysis. Thus, the updated framework is able to perform both stand-alone 
$\nu_{\mu}$-disappearance and $\nu_{e}$-appearance analyses and a complete 3-flavour joint oscillation analysis.
A validation of the new $\nu_{e}$-appearance analysis implemented, compared with the official results, 
is presented in Appendix \ref{app:standalone}.

Neutrino oscillation probabilities are computed for this analysis in a 3-flavour
framework including constant-density matter effects developed by the {\em VALOR} group.
Details of the oscillation probability calculations, estimates of the numerical accuracy of such
calculations and comparisons with other oscillation probability codes 
are given in Appendices A and B of Ref.\ \cite{OurRun123DisappAnalysis3Flavour}. The Fogli-Lisi convention 
is used for the atmospheric mass-squared splitting, as explained in Sec.\ref{sec:osc_params_choice}.

A total of 64 efficiency, energy scale, neutrino flux, neutrino cross-section, 
initial-state nuclear environment and final-state re-scattering systematics are considered for this analysis.
These systematic parameters and their effects 
on the extrapolated single $\mu$-like ring event and single e-like ring event reconstructed 
energy spectrum are described in Sec. \ref{sec:systematics}.

As in previous analyses performed by the {\em VALOR} group, 
the oscillation fit was made using the likelihood-ratio method \cite{StatAnalysisCowan}.
All 64 systematics \nolinebreak considered in this analysis were allowed to float in the fit.
The oscillation fitting method is presented in detail in Sec.\ \ref{sec:neutrino_oscillation_fitting}, including the description 
of the chosen goodness-of-fit test, suitable for the Run 1+2+4 statistics and already tested for the Run 1+2+3 dataset in Ref.
\cite{OurJointAnalysis}.


The fitter performance has been tested in detail using toy MC experiments 
and the results of these tests are presented Sec. \ref{sec:validation}.
Joint 3-flavour oscillation fits of fake datasets have been studied in detail and the results are presented in Sec. \ref{sec:fake_datasets}. These fits were performed using the complete {\em VALOR} analysis with an integrated 
exposure of 6.5700 $\times$10$^{20}$ POT. 


The results of the joint measurement of $|\Delta m^{2}|$, sin$^{2}\theta_{23}$, sin$^{2}\theta_{13}$ and
$\delta_{CP}$, and the systematic parameter pulls, using the Run 1+2+3+4 dataset are given in Sec.\ \ref{sec:results_bestfit_params}, 
and the goodness-of-fit tests are presented in Sec.\ \ref{sec:results_gof}.
The best-fit single $\mu$-like ring and single e-like ring event reconstructed energy spectra 
are shown in Sec.\ \ref{sec:results_bestfit_spectra}.
In Sec.\ \ref{sec:results_confidence_regions}, profile $\Delta\chi^{2}$ distributions are presented 
for each of $|\Delta m^{2}|$, sin$^{2}\theta_{23}$, sin$^{2}\theta_{13}$, $\delta_{CP}$
and 68\% and 90\% confidence regions are presented in
following four 2-dimensional oscillation parameter spaces:
($\sin^{2}\theta_{23}$, $|\Delta m^{2}|$),
($\sin^{2}\theta_{13}$, $|\Delta m^{2}|$),
($\sin^{2}\theta_{13}$, $\delta_{CP}$) and
($\sin^{2}\theta_{23}$, $\sin^{2}\theta_{13}$).
Comparisons of the Run 1+2+3+4 joint 3-flavour oscillation fit results in both normal and inverted hierarchies were made. 
In Sec.\ \ref{sec:results_reactor} results are presented for the Run1+2+3+4 dataset fit including the reactor constraint: 
profile $\Delta\chi^{2}$ distributions are presented 
for each of $|\Delta m^{2}|$, sin$^{2}\theta_{23}$, sin$^{2}\theta_{13}$, $\delta_{CP}$
and 68\% and 90\% confidence regions are presented in
following six 2-dimensional oscillation parameter spaces:
($\sin^{2}\theta_{23}$, $|\Delta m^{2}|$),
($\sin^{2}\theta_{13}$, $|\Delta m^{2}|$),
($\sin^{2}\theta_{13}$, $\delta_{CP}$),
($\sin^{2}\theta_{23}$, $\sin^{2}\theta_{13}$)
($\sin^{2}\theta_{23}$, $\delta_{CP}$) and 
($|\Delta m^{2}|$, $\delta_{CP}$),
comparison also results in both normal and inverted hierarchies.

Finally, a summary is given in Sec.\ \ref{sec:summary}.  




